Tax authorities must identify from millions of self assessed tax returns, the small number of reports who have anomaly. Confirmed non-compliance labels are rare and selectively observed based on auditors point of view, hence using machine learning to identify those anomaly in taxpayers would save time and effort. Unsupervised anomaly detection could be a solution to this problem, but scoring every taxpayer against the entire population could bring two issues. One is that returns that are irregular, and returns that are rare kind, such as landlords, sole traders or self-funded retirees. This paper proposes and implements a four-stage framework that segments taxpayers into peer groups using K-Means method, then detects anomalous reporting within each group using Isolation Forest, then tests how far anomaly mechanisms can be learned by a supervised Random Forest method, and finally explains both segment and anomaly using SHAP. K-Means selected eight well separated segments with silhouette 0.65 that correspond to recognisable taxpayer types. And for every segment, we run Isolation Forest on 56 behavioural and amount features. Then compared it with a population-wide forest scores. Random Forest and a gradient-boosting has been trained on variables created by injecting five documented anomaly mechanisms. SHAP explains the segmentation through a surrogate model and the forests through exact TreeExplainer attributions. Segmentation removed the population-wide forest's overflagging of rare taxpayer types, but it did not raise recall of the injected anomalies. The Random Forest model recovered almost all injected records with F1 score 0.917 and PR-AUC 0.975. Training to supervised model was done using injected anomalies, not to real non-compliance records. Flagged return in this research is rather a statistical anomaly, not a compliance-risk indicator.

\vspace{0.5em}
\noindent\textbf{Keywords:} tax anomaly detection; taxpayer classification; taxpayer segmentation; unsupervised machine learning; supervised machine learning; Isolation Forest; clustering; K-Means; explainable AI; SHAP; tax compliance; anomaly detection; tax risk profiling.
